\section{Conclusion}
\label{sec:conclusion}

We presented a multi-scale sensor-to-G-code reconstruction mismatch framework for detecting adversarial modifications to CNC machining programs.
An autoregressive decoder reconstructs the G-code token sequence that should have produced the observed multi-modal sensor signals; anomalies manifest as elevated reconstruction error when the claimed G-code is inconsistent with the physical process.

No single decoder configuration or scoring method dominates all attack types.
Short-context decoding (V7) achieves strong structural-attack detection through concentrated per-token probability mass, while multi-line decoding (V8) captures sequential consistency needed for semantic attacks.
Rank-based scoring recovers coordinate-perturbation detection where NLL fails, by exploiting relative probability ordering rather than absolute likelihoods.
The per-token confidence gradient of autoregressive G-code prediction---confident numeric tokens (54\% of positions, per-token AUROC = 0.89) versus low-confidence tokens (per-token AUROC near chance)---is the fundamental factor governing coordinate-attack detection and explains why multi-line context, which raises per-token confidence, improves it.

The system is robust to sensor dropout (worst single-board-out AUROC = 0.789), stable across feature configurations (56--110 channels, AUROC variation $< 0.02$), and deployable at low cost (\$492 sensor hardware; measured decoder inference 5.8~ms/window on GPU, 22.3~ms on CPU).
Several negative results constrain the design space: multi-head disagreement scoring provides no signal (AUROC $\approx$ 0.500), embedding-distance methods are blind to token-level attacks, temperature-scaling calibration degrades ECE, and CUSUM achieves only 65.9\% detection at a 38.6\% false alarm rate.

Future work should address five priorities:
\begin{enumerate}
    \item \textbf{Real-time streaming deployment} with position-normalized scoring and CNC controller integration for automatic process halt.
    \item \textbf{Multi-machine and multi-material generalization} on production-grade machines cutting metals, using transfer learning to reduce retraining costs.
    \item \textbf{Attack diagnosis} beyond binary detection, tracing anomaly scores to specific token positions with estimated modification type and magnitude.
    \item \textbf{Defenses against adaptive adversaries.} Our first-order adaptive attacker already nearly evades single-score detection (AUROC $\to$ 0.54); robust detection will require score ensembling, randomized thresholds, or detectors that bound the worst-case perturbation rather than the average case.
    \item \textbf{Sensor integrity assurance} via hardware attestation and cryptographic authentication to prevent sensor spoofing.
\end{enumerate}

The reconstruction-mismatch paradigm---learning the inverse mapping from physical observations to commands and flagging discrepancies---generalizes beyond CNC machining to any cyber-physical system where a digital program controls a physical process (3D printers, robotic welders, semiconductor fabrication, chemical process controllers).
The key requirement is a sensor suite capturing sufficient physical information to constrain command-sequence reconstruction, a condition increasingly met as Industry~4.0 instrumentation becomes standard.

\section*{Reproducibility}
All training scripts, evaluation code, and experiment configurations will be made available upon publication.
The dataset consists of sensor recordings from a Bantam Tools Desktop CNC milling machine; data sharing is subject to institutional approval.
All reported results use 5-fold cross-validation with file-level stratified splits to prevent data leakage.
Attacks are generated deterministically (seed = 42).
Per-fold results are provided in the supplementary material.
